\section{Conclusion}
This controlled extension of \textcite{chen_huang_tien_2021} documents the loss of estimation accuracy and inferential informativeness as excluded-instrument relevance weakens in a high-dimensional IVQR design. Independent compensating noise keeps the marginal treatment distribution and the relationship generating endogeneity fixed. The deterioration affects Oracle-GMM as well as Full-GMM and DML-IVQR, so high-dimensional nuisance estimation cannot be its sole explanation.

Regularized nuisance estimation often improves point accuracy relative to Full-GMM, with exceptions. Coverage, however, is distinct from accuracy and informativeness: near-nominal truth retention can coexist with acceptance of most of the computational domain and little power against false effects. Strong-relevance undercoverage further prevents reading the nominal confidence level as a finite-sample guarantee. Orthogonalization and regularization can mitigate high-dimensional estimation difficulties, but cannot replace identifying information lost when instrument relevance deteriorates.

The findings concern one DGP family, two sample sizes, five quantiles, and three preserved implementations. Density-weight regularity, residualization differences, same-sample estimation, penalty randomness, asymptotic critical values, and finite domains qualify the link to theoretical DML-IVQR. Future controlled comparisons could vary these features individually, examine a DGP satisfying the density conditions, and add sample sizes with systematic domain checks. Such experiments would assess mechanisms without presuming that a particular change restores calibration; no finite-grid diagnostic alone can establish global boundedness.
